\section{经济犯罪 - 生产、销售伪劣商品罪、走私犯罪}
\begin{enumerate}
  \item 很细碎、记忆性的考点。
  \item 生产、销售伪劣产品。
        \begin{enumerate}
          \item 既遂。5W。
          \item 未遂。Opt(销售金额 x 3) + 货值（库存）金额额 \dydy 15W。
          \item 罪数关系。\zd{一般罪名，兜底罪名}。
                \begin{enumerate}
                  \item 定不了劣药罪，那选用此罪。
                  \item \al{卖假烟假酒 \yjt 想象竞合择一重}。
                  \item 侵犯知识产权的，择一重。
                \end{enumerate}
        \end{enumerate}
  \item \kd{食品犯罪}。
        \begin{enumerate}
          \item 生产、销售有毒、有害食品。抽象危险。
                \begin{enumerate}
                  \item \zd{加入有毒、有害的非食品原料}。\al{地沟油、工业酒精、罂粟壳粉末}。
                  \item 食品扩大解释。生的、熟的、半成品、未/加工初级农产品等。\al{用瘦肉精喂猪}。
                \end{enumerate}
          \item 生产、销售不符合安全标准的食品。具体危险。\zd{没有达到毒害性}。
                \begin{enumerate}
                  \item \al{容器被污染、超过保质期、老坛酸菜、过期食品、农药残留、重金属超标}。
                  \item 要求检验检疫、存在国标的。\al{“死”的家禽等}。
                \end{enumerate}
        \end{enumerate}
  \item 药品犯罪。
        \begin{enumerate}
          \item 生产、销售、提供\kd{假药}。抽象危险。
                \begin{enumerate}
                  \item 假药的认定。许可标准 \yjt 实质标准（\zd{成分不符、变质、冒充}）。\zd{吃了没用}。
                  \item \al{用面粉制作速效救心丸 \dui、用面粉 + 白糖制作避孕药 \dui、质量合格、有疗效 \cuo}。
                  \item \al{未经许可、无批号 \cuo \yjt 不看许可了}。
                \end{enumerate}
          \item 生产、销售、提供\kd{劣药}。实害犯（\zd{吃死人}）。
          \item \sjt 提供。
                \begin{enumerate}
                  \item 无偿 \hz 有偿，和销售部分重叠。
                  \item 身份犯。药品\zd{使用}单位及其人员。\al{医院、医生；药厂 \cuo \yjt 生产单位}。
                \end{enumerate}
          \item 结合诈骗。\al{甲购买一批假药，打算免费赠送给老年人，骗取信任，以此诈骗其钱财。
                  然而，存储假药的仓库发生火宅，假药全部被毁。
                  甲构成诈骗罪的犯罪预备 \dui}。
          \item \kd{妨害药品管理}。新增。
                \begin{enumerate}
                  \item \zd{具体危险}。\zd{违禁药、批文、虚假材料申请许可、编造生产/检验记录}。
                  \item \al{《我不是药神》既不构成销售假药、也不构成妨害药品管理 \dui}。
                \end{enumerate}
          \item 非法经营。麻醉、精神药品。
                \begin{enumerate}
                  \item 原料、辅料；药品经营许可证\zd{非法渠道}采购、销售。如\zd{不能认定假药、劣药 \yjt 行政处罚}。
                  \item 麻醉药。\yi 医疗目的贩卖 \yjt 非法经营。\er \zd{牟利目的 \yjt 贩卖毒品}。
                \end{enumerate}
          \item 药品监管渎职。故意、过失均可。
        \end{enumerate}
  \item \sjt 如果行为\zd{同时构成} 8 个具体罪名和生产、销售伪劣产品，\zd{按照处罚较重的论处}。
  \item \kd{走私犯罪}。
        \begin{enumerate}
          \item \xjt 一般法条和特别法条关系。\zd{台阶}。\zd{营利目的}（能否实现，在所不问）。\al{没有卖掉 \dui}。
          \item 走私普通货物、物品。\zd{国家允许}进、出口。逃\zd{关税}。
          \item 走私禁止进出口的货物、物品（兜底罪名）。侵犯\zd{海关管理}制度。\xjt 具体罪名。
                \begin{enumerate}
                  \item 进、出都构成。\zd{假币核武、动物淫秽}。
                        \begin{enumerate}
                          \item 走私武器、弹药；走私核材料；走私假币；走私珍贵动物、珍贵动物制品；走私淫秽物品。
                        \end{enumerate}
                  \item 出口构成。走私\zd{文}物、走私\zd{贵金}属。好的不能出去。
                  \item 进口构成。走私废物。坏的不能进来。
                  \item 除此之外用兜底罪名。\al{走私国家禁止进出口的旧机动车 \yjt 走私禁止进出口的货物、物品}。
                  \item 物品性质定性。能用的正常定。不能用 \yjt 走私普通物品。不能用 + 废 \yjt 走私废物。
                  \item 考试考包容评价，和认识错误结合。\zd{想 A 为 B \yjt A 未遂，B 既遂 \yjt 看 B 的性质}。
                        % https://q.uiver.app/?q=WzAsMyxbMCwwLCLotbDnp4HlhbfkvZPnmoRcXGJsaCJdLFsxLDEsIui1sOengeemgeatoui/m+WHuuWPo+eahFxcYmxoIl0sWzIsMiwi6LWw56eB5pmu6YCaXFxibGgiXSxbMCwxXSxbMSwyXSxbMCwyLCIiLDIseyJjdXJ2ZSI6LTV9XV0=
\[
  \begin{tikzcd}[sep=tiny]
    \txt{走私具体的\blh} \\
    & \txt{走私禁止进出口的\blh} \\
    && \txt{走私普通\blh}
    \arrow[from=1-1, to=2-2]
    \arrow[from=2-2, to=3-3]
    \arrow[curve={height=-30pt}, from=1-1, to=3-3]
  \end{tikzcd}
\]
                        \begin{enumerate}
                          \item \al{甲以为走私的是假币，实际上走私的是淫秽物品}。
                          \item \al{甲让乙帮忙走私淫秽光盘，实际上是废弃的碳渣}。
                          \item \al{甲让乙帮忙走古典音乐光盘，实际上是淫秽光盘}。
                        \end{enumerate}
                \end{enumerate}
        \end{enumerate}
\end{enumerate}
